Swing-up of an underactuated pendulum or cart-pole is a standard test of nonlinear control~\cite{boubaker2012inverted}. We consider two families of solutions. Energy-shaping controllers exploit the physics: they pump energy until the upright value is reached and then hand over to a local stabiliser~\cite{strm2000swinging}. Optimisation-based controllers, such as iLQR/DDP in a receding horizon~\cite{tassa2012synthesis}, are commonly run with a quadratic cost on the state and the control (the algorithm itself is not restricted to it); we hypothesise that such a cost gives a weak signal far from the goal (the far-start success metric tests this).

This paper tests a simple bridge between the two: a \emph{physics-based running cost} $\tfrac12 w_E (E-E^{*})^2$ on the deviation of the mechanical energy from its upright value, used inside iLQR-MPC alone or in addition to the quadratic cost. We ask four questions, fixed before the experiments: (H1) does it raise swing-up success over the quadratic cost, (H2) does it lower control effort, (H3) is it less sensitive to its weight, and (H4) does it match energy shaping with LQR catch?

Contributions are modest and empirical: a controlled comparison of three running costs and an energy-shaping baseline, all reimplemented here, on two tasks with random initial states, one-dimensional weight sweeps, and a comparison in which each cost uses its best swept weight. We report negative and tied outcomes as such (Sec.~\ref{sec:results}).
